\documentclass[border=8pt]{standalone}
\usepackage[UTF8]{ctex}
\usepackage{tikz}
\usepackage{amsmath,amssymb}
\usetikzlibrary{arrows.meta,calc,decorations.markings}
\definecolor{cBlue}{RGB}{37,99,235}
\definecolor{cOrange}{RGB}{234,88,12}
\definecolor{cGray}{RGB}{100,116,139}
\definecolor{cGreen}{RGB}{22,163,74}
\definecolor{cRed}{RGB}{220,38,38}
\begin{document}
\begin{tikzpicture}[scale=1.05,>=Stealth,line width=0.9pt]
  % helix = R (the Lie algebra / the straight line of "time")
  \draw[cRed,line width=1.1pt,domain=0:1440,samples=360,variable=\t]
    plot ({1.7*cos(\t)},{0.45*sin(\t)+\t/1440*3.3});
  \node[cRed,align=left] at (3.0,3.25) {$\mathbb{R}$ （参数线 $t$）\\\footnotesize（即 $\mathfrak{u}(1)=i\mathbb{R}$）};
  % the points above the identity, at t = 0,2pi,4pi,6pi,8pi
  \foreach \z in {0,0.825,1.65,2.475,3.3}{\fill[cRed] (1.7,\z) circle (0.045);}
  \node[cRed,right] at (1.78,0.825) {$2\pi$};
  \node[cRed,right] at (1.78,1.65) {$4\pi$};
  \node[cRed,right] at (1.78,2.475){$6\pi$};
  % base circle = U(1)
  \draw[cBlue,line width=1.4pt] (0,-0.6) ellipse (1.7 and 0.45);
  \node[cBlue] at (0,-1.35) {$U(1)=\mathbb{R}/2\pi\mathbb{Z}$};
  \fill[cGray] (1.7,-0.6) circle (0.05);
  \node[cGray,right] at (1.78,-0.62) {$1$};
  % exp projection: everything above the identity collapses onto 1
  \draw[cGreen,line width=1.1pt,->] (1.7,3.3) .. controls (3.2,1.4) and (3.2,0.2) .. (1.72,-0.55);
  \node[cGreen,align=center] at (3.45,1.0) {$\exp$\\\footnotesize $t\mapsto e^{it}$};
  \draw[cGreen,dashed,->] (1.7,0.825)--(1.7,-0.52);
  \draw[cGreen,dashed,->] (1.7,1.65)--(1.7,-0.5);
  \node[align=center] at (0,-2.05)
    {\footnotesize\color{cGray} $\exp$ 把平直的李代数“卷”成弯曲的群；$0,2\pi,4\pi,\dots$ 全被卷回单位元 $1$。\\[-1pt]
     \footnotesize\color{cGray} 这条直线其实是圆的“万有覆盖”——稍后 SU(2) 盖住 SO(3) 用的是同一套几何。};
\end{tikzpicture}
\end{document}
